\subsection{Induktion und Zerlegung aus den Wortaxiomen}

Die folgenden Hilfssätze machen die Typisierung und die Behandlung eines geordneten Paars als Funktionsargument ausdrücklich verfügbar.

\noindent\begin{minipage}{\linewidth}
\FormulaThmDeltaK[Typisierung eines Anfügungsschritts]{%
\mathsf{WortStr}_A(W,e,s)\dsep u\in W\dsep a\in A\vdash s(u,a)\in W
}{WordStructureStepTyping}{\DeltaRow{Mengen}{A\dsep W\dsep e\dsep u\dsep a}\DeltaRow{Funktionen}{s}}
\end{minipage}\par
\begin{tabproofwide}
  \proofstepwidestar[1]{\mathsf{WortStr}_A(W,e,s)}{\rA}
  \proofstepwidestar[2]{u\in W}{\rA}
  \proofstepwidestar[3]{a\in A}{\rA}
  \proofstepwidestar[1]{e\in W\land s\colon W\times A\to W}{\FormulaRefAuto{WordStructureTypingAxiom}[ax]{1}}
  \proofstepwidestar[2,3]{(u,a)\in W\times A}{\FormulaRefAuto{(a,b)\in A\times B\eqvdash a\in A\land b\in B}[thm]{\rAI{2,3}}}
  \proofstepwidestar[1,2,3]{s(u,a)\in W}{\FormulaRefAuto{F\colon A\to B\dsep x\in A\vdash F(x)\in B}[thm]{\rAEb{4},5}}
\end{tabproofwide}


\noindent\begin{minipage}{\linewidth}
\FormulaThmDeltaK[Injektivität auf geordneten Paaren]{%
\mathsf{WortStr}_A(W,e,s)\dsep p,q\in W\times A\dsep s(p)=s(q)\vdash p=q
}{WordStructurePairInjective}{\DeltaRow{Mengen}{A\dsep W\dsep e\dsep p\dsep q}\DeltaRow{Funktionen}{s}\DeltaRow{Projektionen}{\pi_1\dsep\pi_2}}
\end{minipage}\par
\begin{tabproofwide}
  \proofstepwidestar[1]{\mathsf{WortStr}_A(W,e,s)}{\rA}
  \proofstepwidestar[2]{p,q\in W\times A}{\rA}
  \proofstepwidestar[3]{s(p)=s(q)}{\rA}
  \proofstepwidestar[2]{\pi_1(p)\in W}{\FormulaRefAuto{z\in A\times B\vdash\pi_1(z)\in A}[thm]{\rAEa{2}}}
  \proofstepwidestar[2]{\pi_2(p)\in A}{\FormulaRefAuto{z\in A\times B\vdash\pi_2(z)\in B}[thm]{\rAEa{2}}}
  \proofstepwidestar[2]{\pi_1(q)\in W}{\FormulaRefAuto{z\in A\times B\vdash\pi_1(z)\in A}[thm]{\rAEb{2}}}
  \proofstepwidestar[2]{\pi_2(q)\in A}{\FormulaRefAuto{z\in A\times B\vdash\pi_2(z)\in B}[thm]{\rAEb{2}}}
  \proofstepwidestar[2]{p=(\pi_1(p),\pi_2(p))}{\FormulaRefAuto{z\in A\times B\vdash z=\bigl(\pi_1(z),\pi_2(z)\bigr)}[thm]{\rAEa{2}}}
  \proofstepwidestar[2]{q=(\pi_1(q),\pi_2(q))}{\FormulaRefAuto{z\in A\times B\vdash z=\bigl(\pi_1(z),\pi_2(z)\bigr)}[thm]{\rAEb{2}}}
  \proofstepwidestar[2,3]{s(\pi_1(p),\pi_2(p))=s(\pi_1(q),\pi_2(q))}{\rIE{8,9,3}}
  \proofstepwidestar[1,2,3]{\pi_1(p)=\pi_1(q)\land\pi_2(p)=\pi_2(q)}{\FormulaRefAuto{WordStructureInjectivityAxiom}[ax]{1,4,6,5,7,10}}
  \proofstepwidestar[1,2,3]{p=q}{\rIE{\FormulaRefAuto{a=b\vdash b=a}[thm]{9},\rIE{\rAEa{11},\rAEb{11},8}}}
\end{tabproofwide}





\noindent\begin{minipage}{\linewidth}
\FormulaThmDeltaK[Induktion in einer Wortstruktur]{%
\begin{aligned}[t]&\mathsf{WortStr}_A(W,e,s)\dsep P(e)\dsep{}\\[-2pt]&\forall u\in W\,\forall a\in A\,(P(u)\rightarrow P(s(u,a)))\vdash\forall w\in W\,P(w)\end{aligned}
}{WordStructureInduction}{\DeltaRow{Mengen}{A\dsep W\dsep e\dsep u\dsep a\dsep w}\DeltaRow{Funktionen}{s}\DeltaRow{Prädikat}{P}\DeltaRow{Abkürzung}{U=\{w\in W\mid P(w)\}}}
\end{minipage}\par
\begin{tabproofwide}
  \proofstepwidestar[1]{\mathsf{WortStr}_A(W,e,s)}{\rA}
  \proofstepwidestar[2]{P(e)}{\rA}
  \proofstepwidestar[3]{\forall u\in W\,\forall a\in A\,(P(u)\rightarrow P(s(u,a)))}{\rA}
  \proofstepwidestar[]{U\subseteq W}{\FormulaRefAuto{\{x\in A\mid P(x)\}\subseteq A}[thm]}
  \proofstepwidestar[1]{e\in W\land s\colon W\times A\to W}{\FormulaRefAuto{WordStructureTypingAxiom}[ax]{1}}
  \proofstepwidestar[1,2]{e\in U}{\FormulaRefAuto{x\in\{u\in A\mid P(u)\}\eqvdash x\in A\land P(x)}[thm]{\rAI{\rAEa{5},2}}}
  \proofstepwidestar[7]{u\in U}{\rA}
  \proofstepwidestar[8]{a\in A}{\rA}
  \proofstepwidestar[7]{u\in W\land P(u)}{\FormulaRefAuto{x\in\{u\in A\mid P(u)\}\eqvdash x\in A\land P(x)}[thm]{7}}
  \proofstepwidestar[3,7,8]{P(u)\rightarrow P(s(u,a))}{\rRE{8,\rUE{\rRE{\rAEa{9},\rUE{3}}}}}
  \proofstepwidestar[3,7,8]{P(s(u,a))}{\rRE{\rAEb{9},10}}
  \proofstepwidestar[1,7,8]{s(u,a)\in W}{\FormulaRefAuto{WordStructureStepTyping}[thm]{1,\rAEa{9},8}}
  \proofstepwidestar[1,3,7,8]{s(u,a)\in U}{\FormulaRefAuto{x\in\{u\in A\mid P(u)\}\eqvdash x\in A\land P(x)}[thm]{\rAI{12,11}}}
  \proofstepwidestar[1,3]{\forall u\in U\,\forall a\in A\,s(u,a)\in U}{\rUI{\rRI{7,\rUI{\rRI{8,13}}}}}
  \proofstepwidestar[1,2,3]{U=W}{\FormulaRefAuto{WordStructureMinimalityAxiom}[ax]{1,4,6,14}}
  \proofstepwidestar[16]{w\in W}{\rA}
  \proofstepwidestar[1,2,3,16]{w\in U}{\rIE{\FormulaRefAuto{a=b\vdash b=a}[thm]{15},16}}
  \proofstepwidestar[1,2,3,16]{w\in W\land P(w)}{\FormulaRefAuto{x\in\{u\in A\mid P(u)\}\eqvdash x\in A\land P(x)}[thm]{17}}
  \proofstepwidestar[1,2,3]{\forall w\in W\,P(w)}{\rUI{\rRI{16,\rAEb{18}}}}
\end{tabproofwide}


Die Mengenform und die Prädikatsform der Induktion sind gleichwertig: In der Rückrichtung setzt man \(P(w)\coloneqq w\in U\) und benutzt die Teilmengendefinition.


\Needspace{32\baselineskip}
\noindent\begin{minipage}{\linewidth}
\FormulaThmDeltaK[Induktionsregel in einer Wortstruktur]{%
  \begin{aligned}[t]
    &\mathsf{WortStr}_A(W,e,s)\dsep P(e)\dsep{}\\[-2pt]
    &[u\in W,a\in A,P(u)]\vdots P(s(u,a))\dsep{}\\[-2pt]
    &v\in W\vdash P(v)
  \end{aligned}
}{WordStructureInductionRule}{%
  \DeltaRow{Mengen}{A\dsep W\dsep e\dsep u\dsep a\dsep v}
  \DeltaRow{Funktionen}{s}
  \DeltaRow{Einstelliges Prädikat}{P}
}
\end{minipage}\par
Diese Regel wird aus der durch W3 begründeten Induktion abgeleitet und
gilt für jede Wortstruktur. Wie bei der konkreten Wortregel sind
\(u,a\) beliebig und kommen in den übrigen offenen Annahmen nicht frei
vor. Die Notation \([i,j,k]\vdots l\) in einem Regelverweis nennt die
entlassenen lokalen Annahmen und die Endzeile des Schrittbeweises.
\Needspace{9\baselineskip}
\begin{tabproofwide}
  \proofstepwidestar[1]{\mathsf{WortStr}_A(W,e,s)}{\rA}
  \proofstepwidestar[2]{P(e)}{\rA}
  \proofstepwidestar[3]{v\in W}{\rA}
  \proofstepwidestar[4]{u\in W}{\rA}
  \proofstepwidestar[5]{a\in A}{\rA}
  \proofstepwidestar[6]{P(u)}{\rA}
  \proofstepwidestar[4,5,6]{P(s(u,a))}{%
    \begin{aligned}[t]
      &\text{Schrittprämisse}\\[-2pt]
      &(4,5,6)
    \end{aligned}}
  \proofstepwidestar[]{\forall u\in W\,\forall a\in A\,
    (P(u)\rightarrow P(s(u,a)))}{%
    \rUI{\rRI{4,\rUI{\rRI{5,\rRI{6,7}}}}}}
  \proofstepwidestar[1,2,3]{P(v)}{%
    \rRE{3,\rUE{\FormulaRefAuto{WordStructureInduction}[thm]{1,2,8}}}}
\end{tabproofwide}

\Needspace{16\baselineskip}
\noindent\begin{minipage}{\linewidth}
\FormulaThmDeltaK[Eindeutige Endzerlegung in einer Wortstruktur]{%
\mathsf{WortStr}_A(W,e,s)\dsep w\in W\vdash w=e\lor\exists!p\in W\times A\;w=s(p)
}{WordStructureDecomposition}{\DeltaRow{Mengen}{A\dsep W\dsep e\dsep w\dsep u\dsep a\dsep p\dsep q}\DeltaRow{Funktionen}{s}\DeltaRow{Abkürzung}{P(w)\coloneqq w=e\lor\exists p\in W\times A\;w=s(p)}}
\end{minipage}\par
\begin{tabproofwide}
  \proofstepwidestar[1]{\mathsf{WortStr}_A(W,e,s)}{\rA}
  \proofstepwidestar[2]{w\in W}{\rA}
  \proofstepwidestar[]{e=e}{\rII}
  \proofstepwidestar[]{P(e)}{\rOIa{3}}
  \proofstepwidestar[5]{u\in W}{\rA}
  \proofstepwidestar[6]{a\in A}{\rA}
  \proofstepwidestar[5,6]{(u,a)\in W\times A}{\FormulaRefAuto{(a,b)\in A\times B\eqvdash a\in A\land b\in B}[thm]{\rAI{5,6}}}
  \proofstepwidestar[]{s(u,a)=s(u,a)}{\rII}
  \proofstepwidestar[5,6]{\exists p\in W\times A\;s(u,a)=s(p)}{\rEI{\rAI{7,8}}}
  \proofstepwidestar[5,6]{P(s(u,a))}{\rOIb{9}}
  \proofstepwidestar[1,2]{w=e\lor\exists p\in W\times A\;w=s(p)}{%
    \FormulaRefAuto{WordStructureInductionRule}[thm]{1,4,[5,6]\vdots10,2}}
  \proofstepwidestar[12]{w=e}{\rA}
  \proofstepwidestar[12]{w=e\lor\exists!p\in W\times A\;w=s(p)}{\rOIa{12}}
  \proofstepwidestar[14]{\exists p\in W\times A\;w=s(p)}{\rA}
  \proofstepwidestar[15]{p\in W\times A\land w=s(p)}{\rA}
  \proofstepwidestar[16]{q\in W\times A\land w=s(q)}{\rA}
  \proofstepwidestar[15,16]{s(q)=s(p)}{\FormulaRefAuto{a=b\dsep a=c\vdash b=c}[thm]{\rAEb{16},\rAEb{15}}}
  \proofstepwidestar[1,15,16]{q=p}{\FormulaRefAuto{WordStructurePairInjective}[thm]{1,\rAI{\rAEa{16},\rAEa{15}},17}}
  \proofstepwidestar[1,15]{\forall q\,((q\in W\times A\land w=s(q))\rightarrow q=p)}{\rUI{\rRI{16,18}}}
  \proofstepwidestar[1,15]{\exists!p\in W\times A\;w=s(p)}{\FormulaRefAuto{WordUniqueExistenceFromWitness}[thm]{15,19}}
  \proofstepwidestar[1,15]{w=e\lor\exists!p\in W\times A\;w=s(p)}{\rOIb{20}}
  \proofstepwidestar[1,14]{w=e\lor\exists!p\in W\times A\;w=s(p)}{\rEE{14,15,21}}
  \proofstepwidestar[1,2]{w=e\lor\exists!p\in W\times A\;w=s(p)}{\rOE{11,12,13,14,22}}
\end{tabproofwide}


Die beiden Fälle schließen einander nach \FormulaRefAuto{WordStructureZeroAxiom}[ax] aus. Für ein einelementiges Alphabet ist die Anfügung ein einzelner Nachfolger; die Bedingungen W0 bis W3 entsprechen damit Anfangselement, Nachfolgerabbildung, Injektivität und Induktion der Peano-Struktur.
